\bigskip\Needspace{10\baselineskip}
\pdfbookmark[1]{Pallas Olbersiana. December 1802}{pallas180212}
\begin{center}
PALLAS OLBERSIANA.\par
\medskip\hrule\medskip
Monatliche Correspondenz. Vol. VI. pp.\ 579—583. December 1802.\par
\medskip\hrule
\end{center}\medskip
—\quad —\quad —

Dr.\ \textsc{Gauss} has now compared all these observations by \textsc{Oriani} in Milan with his elements III for the orbit of this planet [p.\ 218 of this volume] and found the following agreement.
\Needspace{22\baselineskip}
\begin{center}\textit{Comparison of \textsc{Oriani}’s observations made with the Milan equatorial sector and Dr.\ \textsc{Gauss}’s elements III for the orbit.}\end{center}
\begin{center}
\begingroup\footnotesize\setlength{\tabcolsep}{2pt}
\begin{tabular}{l|r|r|r|r|r}
1802&\shortstack{Mean time\\in Milan}&\shortstack{Calculated\\RA of Pallas}&\shortstack{Difference\\R—B}&\shortstack{Calculated north\\decl. of $\PallasSymbol$}&\shortstack{Difference\\R—B}\\\hline
10 July&$9^{\mathrm u}\ 40^{\mathrm m}\ 23^{\mathrm s}$&$191^0\ 27'\ 33''$&$\left\{\begin{array}{r}-11''\\-35\end{array}\right.$&$17^0\ 56'\ 20''$&$\left\{\begin{array}{r}-4''\\+4\end{array}\right.$\\
18&$9\ 36\ 4$&$193\ 43\ 47$&$-17$&$17\ 2\ 29$&$-38$\\
24&$9\ 32\ 16$&$195\ 30\ 59$&$-29$&$16\ 19\ 59$&$+12$\\
25&$9\ 44\ 12$&$195\ 49\ 25$&$-23$&$16\ 12\ 40$&$+12$\\
28&$9\ 8\ 6$&$196\ 44\ 25$&$\left\{\begin{array}{r}-33\\-28\end{array}\right.$&$15\ 50\ 56$&$\left\{\begin{array}{r}+1\\+16\end{array}\right.$\\
29&$9\ 11\ 28$&$197\ 3\ 9$&$\left\{\begin{array}{r}-36\\-22\\-36\end{array}\right.$&$15\ 43\ 33$&$\left\{\begin{array}{r}+1\\+18\\+14\end{array}\right.$\\
31&$9\ 6\ 9$&$197\ 40\ 46$&$-34$&$15\ 28\ 55$&$-1$\\
1 Aug.\ &$9\ 4\ 14$&$197\ 59\ 44$&$\left\{\begin{array}{r}-40\\-45\\-49\end{array}\right.$&$15\ 21\ 18$&$\left\{\begin{array}{r}-20\\+2\\+5\end{array}\right.$\\
2&$9\ 23\ 37$&$198\ 19\ 4$&$\left\{\begin{array}{r}-36\\-41\end{array}\right.$&$15\ 13\ 43$&$\left\{\begin{array}{r}+10\\+12\end{array}\right.$
\end{tabular}\endgroup
\end{center}
\SourceLeaf{238}%
\begin{center}
\begingroup\footnotesize\setlength{\tabcolsep}{2pt}
\begin{tabular}{l|r|r|r|r|r}
1802&\shortstack{Mean time\\in Milan}&\shortstack{Calculated\\RA of Pallas}&\shortstack{Difference\\R—B}&\shortstack{Calculated north\\decl. of $\PallasSymbol$}&\shortstack{Difference\\R—B}\\\hline
4 Aug.\ &$8^{\mathrm u}\ 55^{\mathrm m}\ 38^{\mathrm s}$&$198\ 57\ 9$&$\left\{\begin{array}{r}-39''\\-47\end{array}\right.$&$14\ 58\ 52$&$\left\{\begin{array}{r}-5''\\-4\end{array}\right.$\\
5&$9\ 14\ 57$&$199\ 16\ 47$&$-39$&$14\ 51\ 13$&$\left\{\begin{array}{r}+15\\+28\end{array}\right.$\\
6&$9\ 6\ 27$&$199\ 36\ 7$&$\left\{\begin{array}{r}-50\\-51\end{array}\right.$&$14\ 43\ 43$&$-3$\\
7&$9\ 4\ 18$&$199\ 55\ 38$&$\left\{\begin{array}{r}-48\\-47\end{array}\right.$&$14\ 36\ 8$&$+15$\\
8&$9\ 0\ 14$&$200\ 15\ 9$&$-56$&$14\ 28\ 35$&$+14$
\end{tabular}\endgroup
\end{center}

Dr.\ \textsc{Gauss} did later give a fourth set of improved elements (October issue, p.\ 394) [p.\ 223 of this volume], but did not himself compare the Milan observations above with them. He did, however, make the following comparison, from which one can judge roughly how much better elements IV agree with the observations than elements III do.
\begin{center}
Elements IV give\par\medskip
\begin{tabular}{l|r|r}
1802&A greater RA&A smaller declination\\\hline
8 July&$24''$&$8''$\\
4 Aug.\ &$66$&$0$
\end{tabular}\par\medskip
than elements III.
\end{center}

Thus, around the end of the observations, elements III give a right ascension about 1 minute too small, while elements IV give one about $20''$ too large. The declinations agree fairly well in both cases, as far as the accuracy of the observations permits.

Using these latest observations by \textsc{Oriani}, Dr.\ \textsc{Gauss} has now also improved the elements of this orbit for the fifth time and obtained the following quantities.
\Needspace{14\baselineskip}
\begin{center}\textit{Elements V of Pallas Olbersiana.}\end{center}
\begin{center}
\begingroup\small\setlength{\tabcolsep}{3pt}
\begin{tabular}{l r}
\shortstack[l]{Epoch of 31 March 1802\\at noon in Seeberg}&$162^0\ 46'\ 58''2$\\
Aphelion $\ .\ .\ .\ .\ .\ .\ .\ .\ .$&$301\ 28\ 24.0$\\
Ascending node $\ .\ .\ .\ .\ .\ .$&$172\ 27\ 3.0$\\
Mean daily tropical motion $\ .$&$769''583$\\[3pt]
Logarithm of the semimajor axis&$0.4425529$\\
Eccentricity $\ .\ .\ .\ .\ .\ .\ .\ .\ .$&$0.244976$\\
Orbital inclination $\ .\ .\ .\ .\ .\ .\ .$&$34^0\ 37'\ 40''$
\end{tabular}\endgroup
\end{center}

Dr.\ \textsc{Gauss} then compared the latest of all \textsc{Messier}’s observations with these elements (November issue, p.\ 502). He obtained the values below.
\Needspace{12\baselineskip}
\begin{center}
\begingroup\small\setlength{\tabcolsep}{4pt}
\begin{tabular}{l|r|r|r}
1802&\shortstack{Mean time\\in Paris}&\shortstack{Calculated RA\\of Pallas}&\shortstack{Difference\\R—B}\\\hline
21 Sept.\ &$7^{\mathrm u}\ 28^{\mathrm m}$&$215^0\ 49'\ 14''$&$+28''$
\end{tabular}\par\medskip
\begin{tabular}{r|r|r|r}
\shortstack{Calculated declination\\of Pallas}&\shortstack{Difference\\R—B}&\shortstack{Distance\\from the $\EarthSymbol$}&Brightness\\\hline
$9^0\ 0'\ 16''\ \mathrm N.$&$+48''$&$3.518$&$0.485$
\end{tabular}\endgroup
\end{center}

The difference between calculation and observation is therefore still moderate. Dr.\ \textsc{Gauss} writes as follows.
“Since this observation by \textsc{Messier} was made in such an unfavorable position (by a rough calculation, Pallas was only 13 degrees above the horizon), I fully agree with you that one can draw little that is reliable from this \textit{single} observation for improving the elements. In any case, the improvement would probably have no very considerable effect on the position of \textit{Pallas} next year. I have recalculated several positions using elements V
\SourceLeaf{239}%
and found that for 4 Feb., 24 March, 11 May and 28 June 1803, the right ascensions were smaller by $3\frac12$ min., $5\frac12$ min., $7\frac12$ min. and 9 min., respectively, and the declinations were all about 1 min. greater than in the ephemeris calculated from elements IV.” (October issue, p.\ 395.) [p.\ 224 of this volume.]

To leave nothing more to be desired here either, Dr.\ \textsc{Gauss} has recalculated the whole ephemeris for the motion of Pallas in 1803 from elements V and extended it by another month. On 3 April 1803, \textit{Pallas} will begin to exceed the brightness it had on 21 September of this year, and it can then be observed at a much greater altitude above the horizon. There can now be little doubt that it will be visible in our more powerful telescopes. The ephemeris follows. [The ephemeris from 4 Feb.\ to 9 Aug.\ 1803, giving right ascension, declination and brightness in the form used so far.]

In the third column, the unit is the brightness the planet would have at a distance of 1 from the Sun and the Earth. On this scale its brightness in 1802 was as follows.
\begin{center}
\begingroup\small
\begin{tabular}{llr r l}
1802&April&4 $\ .\ .$&$0.08997$&[Manuscript.] More precisely $0.09019$\\
&May&16 $\ .\ .$&$0.04740$&\\
&Aug.&10 $\ .\ .$&$0.01455$&\\
&Sept.&21 $\ .\ .$&$0.01030$&
\end{tabular}\par\medskip
\begin{tabular}{llr r}
1801&Jan.&1 $\ .\ .$&$0.03637$\\
&Feb.&11 $\ .\ .$&$0.02409$\\
&Dec.&7 $\ .\ .$&$0.02494$\\
1802&Feb.&26 $\ .\ .$&$0.05633$
\end{tabular}\endgroup
\end{center}
\begin{flushright}[\textsc{von Zach.}]\end{flushright}

